\documentclass[conference]{IEEEtran}
\IEEEoverridecommandlockouts
% The preceding line is only needed to identify funding in the first footnote. If that is unneeded, please comment it out.
\usepackage{cite}
\usepackage{amsmath,amssymb,amsfonts}
\usepackage{algorithmic}
\usepackage{graphicx}
\usepackage{textcomp}
\usepackage{xcolor}
\usepackage{hyperref}
\usepackage[turkish]{babel}
\usepackage[utf8]{inputenc}
\usepackage{booktabs}
\usepackage{float}
\usepackage{lipsum} % Sadece test ve tasarımı görmek için

\def\BibTeX{{\rm B\kern-.05em{\sc i\kern-.025em b}\kern-.08em
    T\kern-.1667em\lower.7ex\hbox{E}\kern-.125emX}}

\begin{document}

\title{Araba Gövde Tipi Sınıflandırma: Derin Öğrenme Yaklaşımları ve Performans Analizi\\
{\Large Kocaeli Üniversitesi Yazlab 2 - 3.Proje}
}

\author{\IEEEauthorblockN{1\textsuperscript{st} Berat Batuhan Şahin}
\IEEEauthorblockA{\textit{Bilgisayar Mühendisliği Bölümü} \\
\textit{Kocaeli Üniversitesi}\\
230201026 \\
Kocaeli, Türkiye}
\and
\IEEEauthorblockN{2\textsuperscript{nd} Emirhan Çalık}
\IEEEauthorblockA{\textit{Bilgisayar Mühendisliği Bölümü} \\
\textit{Kocaeli Üniversitesi}\\
230201108 \\
Kocaeli, Türkiye}
}

\maketitle

\begin{abstract}
Bilgisayarlı görü (computer vision) sistemleri, günümüzde otonom sürüşten akıllı şehir planlamasına kadar pek çok alanda kritik bir rol oynamaktadır. Bu projede, araçların gövde tiplerini görsel veriler üzerinden otomatik olarak sınıflandırabilen derin öğrenme tabanlı bir modelin geliştirilmesi amaçlanmıştır. Proje kapsamında SUV, VAN, Station Wagon, Micro, Açık Tekerlekli, Sedan, Hatchback ve Pick Up olmak üzere 8 farklı araç gövde sınıfı belirlenmiştir. Açık kaynaklı çeşitli veri setlerinden derlenerek oluşturulan özelleştirilmiş veri seti, modelin farklı açılar, ışık koşulları ve arka planlar altındaki genelleme kapasitesini (generalization) artırmak amacıyla kullanılmıştır. Ön işleme aşamasında görüntülerin en-boy oranını bozmayan Letterboxing tekniği, ImageNet standartlarında normalizasyon ve sentetik veri çoğaltma (data augmentation) yöntemleri uygulanmıştır. Model mimarisi olarak, derin ağlarda sıkça karşılaşılan kaybolan gradyan problemini aşmak üzere tasarlanmış ResNet-50 ağı transfer öğrenme (transfer learning) yöntemiyle projeye entegre edilmiştir. Sınıflandırma katmanlarına uygulanan Dropout ve Batch Normalization teknikleri ile aşırı uyum (overfitting) problemi engellenmiştir. Model, eğitim sırasında hiç görmediği test verisi üzerinde \%89.75 Doğruluk (Accuracy) ve 0.9204 Macro F1-Skoru elde etmiştir. Ayrıca, eğitilen model Flask/Streamlit tabanlı bir web arayüzüne entegre edilerek son kullanıcıya hitap eden, gerçek zamanlı çalışan ve olasılık dağılımlarını görselleştiren uçtan uca bir ürün haline getirilmiştir.
\end{abstract}

\begin{IEEEkeywords}
Derin Öğrenme, Evrişimli Sinir Ağları, ResNet-50, Araç Sınıflandırma, Görüntü İşleme, Otonom Sistemler
\end{IEEEkeywords}

\section{Giriş}
Yapay zeka ve özellikle derin öğrenme algoritmalarının gelişmesiyle birlikte, nesne tanıma ve sınıflandırma problemleri klasik görüntü işleme yöntemlerinin ötesine geçmiş, çok daha karmaşık problemler çözülebilir hale gelmiştir. Otomotiv ve ulaşım sektöründe, güvenlik kameralarından, otoyol gişelerinden veya otonom araçların sensörlerinden akan verilerin anlamlandırılması büyük önem taşımaktadır. Bu bağlamda, bir aracın sadece marka veya modelinin değil, doğrudan yapısal özellikleri olan "gövde tipinin" (body type) tespit edilmesi; trafik yoğunluğu analizinden pazar araştırmalarına, hatta suçla mücadele yöntemlerine kadar geniş bir yelpazede kullanılabilmektedir.

Ancak, araç gövde tiplerinin sınıflandırılması, kendi içinde birçok zorluk (challenge) barındıran bir problemdir. Bir SUV ile Station Wagon veya bir Hatchback ile Sedan arasındaki görsel farklılıklar bazen sadece arka tavan çizgisi veya bagaj hacmi ile sınırlı kalmaktadır. Bunun yanı sıra araçların fotoğraflarının çekildiği açılar, günün saatine bağlı olarak değişen aydınlatma koşulları, yansımalar ve arka plandaki karmaşıklıklar, algoritmaların "öznitelik (feature) çıkarımı" yapmasını zorlaştırmaktadır. 

Bu projede, belirtilen bu zorlukların üstesinden gelmek ve yüksek doğruluk oranlarına ulaşmak adına derin öğrenme temelli bir Evrişimli Sinir Ağı (Convolutional Neural Network - CNN) yaklaşımı benimsenmiştir. Literatürde görüntü sınıflandırma denildiğinde akla ilk gelen ve çok derin mimarilerin eğitilmesine olanak tanıyan ResNet (Residual Network) mimarisi tercih edilmiştir. Modelin sadece eğitim setini ezberlemesini değil, gerçek hayat senaryolarına adapte olmasını (genelleme) sağlamak adına çeşitli regularizasyon (düzenlileştirme) adımları uygulanmıştır. Çalışmanın nihai hedefi, sadece bir model eğitmek değil, aynı zamanda eğitilmiş bu modeli gerçek dünyada kullanılabilecek interaktif bir web arayüzü ile sunmaktır.

Makalenin geri kalanı şu şekilde organize edilmiştir: Bölüm II'de veri setinin oluşturulma süreci, karşılaşılan zorluklar ve veri ön işleme adımları detaylandırılmıştır. Bölüm III, modelin mimarisini ve eğitim sürecini anlatmaktadır. Bölüm IV'te deneysel sonuçlar, grafiksel analizler ve metrikler sunulmuş, Bölüm V'te ise geliştirilen web arayüzü ve entegrasyon süreçleri açıklanmıştır. Makale, Bölüm VI'daki genel değerlendirme ve sonuç ile tamamlanmaktadır.

\section{Veri Seti ve Ön İşleme Adımları}

\subsection{Veri Setinin Oluşturulması ve Organizasyonu}
Derin öğrenme modelleri, başarılı bir şekilde eğitilebilmek için büyük miktarda ve çeşitliğe sahip veriye ihtiyaç duyar. İstenen 8 farklı sınıfın (SUV, VAN, Station Wagon, Micro, Açık Tekerlekli, Sedan, Hatchback, Pick Up) her birine ait yeterli miktarda verinin tek bir veri setinde bulunmaması sebebiyle, açık kaynaklı farklı platformlardan çoklu veri setleri incelenmiş ve birleştirilmiştir. Projede referans ve temel veri kaynağı olarak Kaggle platformunda bulunan üç büyük veri setinden yararlanılmıştır:
\begin{enumerate}
    \item \textbf{Cars Body Type Cropped} \cite{ref_dataset1}: Araçların arka planlardan nispeten arındırıldığı ve doğrudan gövdeye odaklanıldığı, temiz etiketlere sahip bir veri setidir.
    \item \textbf{Stanford Car Body Type Data} \cite{ref_dataset2}: Daha çok akademik amaçlarla toplanmış, farklı açı ve modellerden yüzlerce aracın bulunduğu köklü bir veri setidir.
    \item \textbf{Gigaflexhicle} \cite{ref_dataset3}: Modern araçların ve farklı ışık/kamera koşullarının bulunduğu geniş kapsamlı bir veri havuzudur.
\end{enumerate}

Verilerin birleştirilmesi aşamasında, her sınıfın eşit miktarda örneğe sahip olması için dengeleme (balancing) işlemi yapılmıştır. Bir sınıfın diğerinden çok daha fazla örneğe sahip olması, modelin o sınıfa karşı önyargılı (bias) olmasına sebep olabilmektedir. Eksik kalan sınıflar için veri artırımı uygulanarak sınıflar arası denge güvence altına alınmıştır. Eğitim sürecinde doğrulama seti (Validation) ayırımı yapılarak modelin hiperparametreleri ayarlanmış, ayrıca en son yapılacak test için modelin eğitim sırasında "hiç görmediği" ayrı bir Test seti ayrılmıştır.

\subsection{Görüntü Ön İşleme (Image Preprocessing)}
Görüntülerin modele verilmeden önce belirli standartlara oturtulması zorunludur. Derin öğrenme modelleri genellikle sabit boyutta matris girdilerine ihtiyaç duyarlar.
Ön işleme adımları aşağıda sırasıyla listelenmiştir:

\textbf{1. Boyutlandırma ve Letterbox Tekniği:} 
ResNet-50 mimarisi varsayılan olarak $224 \times 224$ piksel boyutunda girdi bekler. Standart yeniden boyutlandırma (resize) işlemleri, en-boy oranı (aspect ratio) farklı olan dikdörtgen araç fotoğraflarının sıkışmasına veya genişlemesine sebep olarak, modelin fiziksel özellikleri yanlış öğrenmesine yol açar. Bu problemi aşmak için "Letterbox to Square" yöntemi kullanılmıştır. Görüntü orijinal en-boy oranı korunarak küçültülmüş ve $224 \times 224$'lük gri bir tuvalin (canvas) ortasına yerleştirilmiştir. Bu sayede araç silüetinde bozulma yaşanmamıştır.

\textbf{2. Normalizasyon:}
Derin ağların daha hızlı ve stabil yakınsaması (convergence) için giriş değerlerinin standardize edilmesi gerekir. Görüntüler $[0, 1]$ piksel değer aralığına getirildikten sonra, ImageNet veritabanının istatistiksel verileri baz alınarak normalize edilmiştir. Denklem (\ref{eq:norm})'de gösterildiği üzere her bir kanal (R, G, B) ortalamadan (Mean) çıkarılmış ve standart sapmaya (Std) bölünmüştür.
\begin{equation}
    x_{norm} = \frac{x_{orijinal} - \mu}{\sigma} \label{eq:norm}
\end{equation}
Burada $\mu = [0.485, 0.456, 0.406]$ ve $\sigma = [0.229, 0.224, 0.225]$ olarak tanımlanmıştır.

\textbf{3. Veri Çoğaltma (Data Augmentation):}
Modelin ezberlemesini önlemek için veri setini sentetik olarak büyüten augmentasyon işlemleri sadece "Eğitim" (Training) setine uygulanmıştır. RandomHorizontalFlip ile sağ-sol açıları genelleştirilmiş, RandomRotation ile eğimli çekimler simüle edilmiş, ColorJitter ile parlaklık, kontrast ve doygunluk değişiklikleri yapılarak farklı saat/kamera kondisyonları taklit edilmiştir.

\section{Model Mimarisi ve Eğitim Stratejisi}

\subsection{Evrişimli Sinir Ağları ve ResNet-50}
Evrişimli Sinir Ağları (CNN), görüntüdeki lokal özellikleri (kenar, köşe, doku vb.) evrişim (convolution) filtreleriyle kademeli olarak öğrenen derin öğrenme mimarileridir. Geleneksel CNN mimarileri derinleştikçe, geri yayılım (backpropagation) esnasında gradyanların sıfıra yaklaşması olarak bilinen "kaybolan gradyan" (vanishing gradient) problemi ortaya çıkmaktadır. Bu sorun, ResNet (Residual Network) mimarisi \cite{ref_resnet} ile "Skip Connection" (Atlama Bağlantıları) kullanılarak çözülmüştür. 

Bir Residual Block içindeki formülasyon $H(x) = F(x) + x$ şeklindedir. Burada $x$ giriş vektörü, $F(x)$ iki evrişim katmanının öğrendiği fonksiyon, $H(x)$ ise bloğun toplam çıktısıdır. Bu toplama işlemi sayesinde gradyanlar atlama bağlantısı üzerinden kesintisiz akabilmektedir.

\subsection{Transfer Öğrenme ve Karar Katmanı Tasarımı}
Projede ResNet-50 modeli ImageNet verisiyle önceden eğitilmiş ağırlıklarla (pre-trained) projeye dahil edilmiştir. Ağın düşük seviyeli özellikleri (çizgiler, renkler vb.) çıkaran ilk katmanları korunmuş, ağın sonundaki orijinal 1000 sınıflı sınıflandırma katmanı çıkarılarak yerine bizim 8 sınıfımızı temsil eden özel bir katman eklenmiştir.
Özel Tam Bağlantılı (Fully Connected) katman mimarimiz şöyledir:
\begin{itemize}
    \item \textbf{Linear Layer 1:} $2048 \rightarrow 1024$ nöron
    \item \textbf{Batch Normalization} ve \textbf{ReLU} aktivasyonu
    \item \textbf{Dropout (p=0.5):} Nöronların \%50'si rastgele kapatılarak aşırı uyum önlenmiştir.
    \item \textbf{Linear Layer 2:} $1024 \rightarrow 512$ nöron
    \item \textbf{Batch Normalization} ve \textbf{ReLU} aktivasyonu
    \item \textbf{Dropout (p=0.4)}
    \item \textbf{Çıkış Katmanı (Output Layer):} $512 \rightarrow 8$ sınıf
\end{itemize}

\subsection{Optimizasyon ve Kayıp Fonksiyonu}
Çok sınıflı sınıflandırma problemi için model çıkışlarına Softmax aktivasyonu uygulanarak her bir sınıfın olasılık değeri hesaplanmış ve kayıp fonksiyonu olarak Cross-Entropy Loss kullanılmıştır. Cross Entropy formülü denklem (\ref{eq:ce}) ile verilmektedir:
\begin{equation}
    Loss = - \sum_{i=1}^{C} y_i \log(\hat{y}_i) \label{eq:ce}
\end{equation}
Burada $C=8$, $y_i$ gerçek etiket dağılımını, $\hat{y}_i$ ise modelin tahmin ettiği olasılık dağılımını temsil etmektedir.

Model ağırlıklarının güncellenmesinde Adam (Adaptive Moment Estimation) \cite{ref_adam} optimizasyon algoritması kullanılmıştır. Eğitim süresince "ReduceLROnPlateau" mekanizması ile validation loss'un düşmesinin durduğu noktalarda öğrenme oranı (Learning Rate) otomatik olarak düşürülmüştür.

\begin{figure*}[htbp]
\centerline{\includegraphics[width=0.85\textwidth, height=8cm]{loss_graph.png}} % BURAYA loss grafiği gelecek
\caption{Eğitim ve Doğrulama (Validation) Kayıp (Loss) Grafiği. X ekseni epoch sayısını, Y ekseni ise Cross-Entropy kayıp değerini göstermektedir.}
\label{fig:loss}
\end{figure*}

\begin{figure*}[htbp]
\centerline{\includegraphics[width=0.85\textwidth, height=8cm]{acc_graph.png}} % BURAYA accuracy grafiği gelecek
\caption{Eğitim ve Doğrulama Accuracy (Doğruluk) Grafiği. Eğriler, modelin çağlar (epoch) boyunca doğruluğunun nasıl istikrar kazandığını göstermektedir.}
\label{fig:acc}
\end{figure*}

\section{Deneysel Sonuçlar ve Metrikler}

Modelin başarısı, daha önce hiç görmediği bir test veri seti üzerinde sınanmıştır. Tahminler esnasında Test-Time Augmentation (TTA) işlemi kullanılmıştır. Bu işlemde test görüntüsünün orijinal hali ile yatay çevrilmiş (flipped) halinin model tarafından üretilen karar olasılıklarının ortalaması alınmış, böylece sonuçların daha güvenilir ve stabil (robust) olması sağlanmıştır.

\subsection{Değerlendirme Metrikleri}
Sınıflandırma problemlerinde sadece Accuracy (Doğruluk) yeterli bir ölçüt değildir, zira dengesiz veri setlerinde yanıltıcı olabilir. Bu nedenle Precision (Kesinlik), Recall (Duyarlılık) ve F1-Score gibi daha güvenilir metrikler formüle edilmiştir:
\begin{equation}
    Precision = \frac{TP}{TP + FP}
\end{equation}
\begin{equation}
    Recall = \frac{TP}{TP + FN}
\end{equation}
\begin{equation}
    F_1 = 2 \times \frac{Precision \times Recall}{Precision + Recall}
\end{equation}

Modelimiz üzerinde yapılan test sonucunda \textbf{Accuracy (Doğruluk) \%89.75} olarak bulunurken, modelin asıl güvenilirliğini gösteren \textbf{Macro F1-Skoru 0.9204} gibi oldukça başarılı bir seviyeye ulaşmıştır.

Tablo \ref{tab:classification_report}'ta görüldüğü üzere model özellikle "VAN", "MICRO", "AÇIK TEKERLEK" ve "PICK UP" sınıflarında neredeyse hatasız çalışmış ve F1 skorları 0.97 - 0.99 bandına çıkmıştır.

\begin{table}[htbp]
\caption{Sınıf Bazlı Performans Raporu (Classification Report)}
\begin{center}
\begin{tabular}{lcccc}
\toprule
\textbf{Sınıf} & \textbf{Precision} & \textbf{Recall} & \textbf{F1-Score} & \textbf{Support}\\
\midrule
1. SUV           & 0.91 & 0.91 & 0.91 & 235 \\
2. VAN           & 0.98 & 1.00 & 0.99 & 111 \\
3. STATION WAGON & 0.91 & 0.79 & 0.85 & 231 \\
4. MICRO         & 0.99 & 0.99 & 0.99 & 72  \\
5. AÇIK TEKERLEK & 1.00 & 0.97 & 0.99 & 116 \\
6. SEDAN         & 0.85 & 0.88 & 0.87 & 275 \\
7. HATCHBACK     & 0.77 & 0.83 & 0.80 & 226 \\
8. PICK UP       & 0.97 & 0.98 & 0.97 & 159 \\
\midrule
\textbf{Accuracy}  & \multicolumn{3}{c}{\textbf{0.90}} & 1425 \\
\textbf{Macro Avg} & \textbf{0.92} & \textbf{0.92} & \textbf{0.92} & 1425 \\
\bottomrule
\end{tabular}
\label{tab:classification_report}
\end{center}
\end{table}

\subsection{Karışıklık Matrisi (Confusion Matrix) Analizi}

Şekil \ref{fig:cm}'de gösterilmesi planlanan Karışıklık Matrisinin detaylı incelenmesinden elde edilen sonuçlara göre model, görsel olarak birbirine benzeyen sınıfları karıştırmaya daha meyillidir. Örneğin Station Wagon (SW), Sedan ve Hatchback sınıfları genel şasi yapısı itibariyle oldukça benzer formlara sahiptir; model hatalarının (False Positives/False Negatives) büyük bir kısmı bu 3 sınıf arasındaki geçişlerde yaşanmıştır. Hatchback sınıfının F1 skorunun 0.80 seviyesinde kalarak en düşük başarıyı göstermesinin temel sebebi, büyük Hatchback'lerin SUV veya kısa Station Wagon'larla benzerlik göstermesidir. 
Ancak Pick-Up (kasa yapısı) ve Açık Tekerlek (Formula araçları vb.) sınıflarının belirgin geometrik özellikleri sayesinde hata payları minimize edilmiştir.

\begin{figure*}[htbp]
\centerline{\includegraphics[width=0.85\textwidth, height=8cm]{confusion_matrix.png}} % BURAYA cm grafiği gelecek
\caption{Normalize Edilmiş Karışıklık Matrisi (Normalized Confusion Matrix). Köşegen (diagonal) üzerindeki yüksek değerler modelin doğru tahminlerini temsil etmektedir.}
\label{fig:cm}
\end{figure*}

\section{Web Arayüzü Geliştirme ve Entegrasyonu}

Akademik ve teorik ortamda eğitilen bir derin öğrenme modelinin, son kullanıcı tarafından rahatlıkla test edilebilmesi ve kullanışlı bir ürün haline gelmesi için web tabanlı bir Grafiksel Kullanıcı Arayüzü (GUI) geliştirilmiştir.

\subsection{Kullanıcı Deneyimi (UX) ve Arayüz Tasarımı}
Geliştirilen arayüz, kullanım kolaylığını maksimize etmek üzere tasarlanmıştır. Arayüzün ana fonksiyonları şunlardır:
\begin{itemize}
    \item \textbf{Görüntü Yükleme Alanı:} Sürükle ve bırak (Drag \& Drop) destekli, .jpg, .png formatlarındaki görsellerin kolayca sisteme yüklenebildiği alan.
    \item \textbf{Önizleme (Preview):} Yüklenen görselin sisteme aktarılır aktarılmaz kullanıcı tarafından doğrulanabilmesi için ekranda gösterilmesi.
    \item \textbf{Hızlı Sınıflandırma (Inference):} Modelin GPU destekli sunucuda ağırlıkları yükleyerek, yüklenen görüntüyü saniyeler içinde analiz etmesi.
\end{itemize}

\subsection{Tahmin Çıktısı ve Olasılık Dağılım Grafikleri}
Kullanıcı "Tahmin Et" butonuna tıkladıktan sonra, sistem sonucu büyük ve belirgin bir etiketle sunmaktadır (Örn: "Tahmin Edilen Sınıf: \textbf{SUV}"). 
Bununla da kalmayıp, arka planda çalışan ResNet-50 modelinin Softmax katmanından elde ettiği 8 farklı olasılık değeri yüzdelik bir çubuk grafiğe (Bar Chart) dönüştürülmektedir. Böylece kullanıcı modelin kararı verirken ikinci veya üçüncü ihtimal olarak hangi gövde tiplerini ne kadar olasılıkla düşündüğünü şeffaf bir şekilde analiz edebilmektedir. 

Ayrıca, arayüzün daha hızlı çalışabilmesi ve sunucu bellek tüketimini azaltmak amacıyla büyük boyutlu Float32 ağırlıkları Float16 (FP16) standardına sıkıştırılarak depolanmıştır. Bu optimizasyon tahmin doğruluk oranlarında kayda değer bir düşüş yaratmazken, yükleme hızlarında avantaj sağlamıştır.

\begin{figure*}[htbp]
\centerline{\includegraphics[width=0.85\textwidth, height=8cm]{gui_screenshot.png}} % BURAYA Arayüz ekran görüntüsü gelecek
\caption{Geliştirilen Web Arayüzünün görünümü. Yüklenen orijinal görsel, tahmin sonucu ve tüm sınıflara ait olasılık dağılım çubuk grafiği (Bar Chart) birlikte gösterilmektedir.}
\label{fig:gui}
\end{figure*}

\section{Sonuç ve Gelecek Çalışmalar}
Otonom sistemler ve akıllı trafik yönetim ağları için kritik bir alt yapı sağlayan araç gövde tipi sınıflandırması, bu projede uçtan uca başarılı bir şekilde gerçekleştirilmiştir. Dört farklı kaynaktan toplanarak oluşturulan veri seti üzerinde, ImageNet ağırlıklarıyla başlatılan ResNet-50 modeli kullanılmıştır.

Uygulanan Letterboxing, veri artırımı (augmentation) ve Dropout/Early Stopping stratejileri ile overfitting problemi büyük oranda çözülmüş olup modelin genelleme başarısı kanıtlanmıştır. Eğitim esnasında hiç kullanılmayan test verisi ile yapılan değerlendirmede \textbf{0.9204 Macro F1} skoru elde edilmiş, modelin zorluklara karşı oldukça dayanıklı (robust) olduğu gözlemlenmiştir. Eğitilen yüksek kapasiteli modelin, bir web arayüzüne başarıyla entegre edilmesi ile proje ürün aşamasına taşınmıştır.

Gelecek çalışmalarda veri seti boyutunun daha da artırılması, özellikle benzer özellikler sergileyen Sedan-Hatchback gibi alt segmentlerdeki ayrımın netleştirilebilmesi için nesne tespiti (object detection) ve dikkat mekanizmalarını (attention mechanisms) temel alan Swin Transformer, Vision Transformer (ViT) gibi modern mimarilerin incelenmesi planlanmaktadır.

\begin{thebibliography}{00}
\bibitem{ref_dataset1} A. Boukhris, ``Cars Body Type Cropped Dataset,'' Kaggle. [Çevrimiçi]. Erişim: https://www.kaggle.com/datasets/ademboukhris/cars-body-type-cropped
\bibitem{ref_dataset2} M. Mahurkar, ``Stanford Car Body Type Data,'' Kaggle. [Çevrimiçi]. Erişim: https://www.kaggle.com/datasets/mayurmahurkar/stanford-car-body-type-data
\bibitem{ref_dataset3} D. Gaus, ``Gigaflexhicle Dataset,'' Kaggle. [Çevrimiçi]. Erişim: https://www.kaggle.com/datasets/dmitrygaus/gigaflexhicle?utm
\bibitem{ref_resnet} K. He, X. Zhang, S. Ren, ve J. Sun, ``Deep Residual Learning for Image Recognition,'' in Proceedings of the IEEE conference on computer vision and pattern recognition, 2016, ss. 770--778.
\bibitem{ref_adam} D. P. Kingma ve J. Ba, ``Adam: A Method for Stochastic Optimization,'' in International Conference on Learning Representations (ICLR), 2015.
\bibitem{ref_pytorch} A. Paszke vd., ``PyTorch: An Imperative Style, High-Performance Deep Learning Library,'' in Advances in Neural Information Processing Systems 32, Curran Associates, Inc., 2019, ss. 8024--8035.
\end{thebibliography}

\end{document}
